\documentclass[11pt,a4paper]{article}

\usepackage[margin=2.3cm]{geometry}
\usepackage{amsmath,amssymb}
\usepackage{graphicx}
\usepackage{booktabs}
\usepackage{array}
\usepackage{xcolor}
\usepackage{enumitem}
\usepackage[hidelinks]{hyperref}
\usepackage{caption}

\emergencystretch=1.5em
\captionsetup{font=small,labelfont=bf,width=.92\textwidth}
\setlength{\parskip}{4pt plus 1pt}
\setlength{\parindent}{0pt}

\newcommand{\code}[1]{\texttt{#1}}
\newcommand{\neff}{n_{\mathrm{eff}}}
\newcommand{\chigeo}{\chi_{\mathrm{geo}}}
\newcommand{\Pisurr}{\Pi}
\newcommand{\EMA}{\operatorname{EMA}}

\title{Order parameter and mechanism of the RBLapSum
\code{through\_rank\_kappa} instability\\[2pt]
\large Research note: one-step score-scale dynamics}
\author{}
\date{2026-09-16/17. Follow-up to the kappa stability investigation.}

\begin{document}
\maketitle
\vspace{-2em}
\tableofcontents

% ============================================================
\section{Purpose and summary of conclusions}
% ============================================================

The previous investigation
(\code{docs/rblapsum-kappa-stability-neutral.tex}) introduced
\[
\chigeo=\frac{b}{2T\delta}
\]
as an empirical instability indicator. This note tests whether \(\chigeo\) is a fundamental control parameter or only a useful proxy, and investigates the mechanism responsible for score-inflation bursts.

Four hypotheses are considered: H1, that \(\chigeo\) is an empirical proxy rather than a fundamental order parameter; H2, that the dimensionless surrogate gain
\[
\Pisurr=\|L_\kappa D_z\|_F^2
\]
is a more appropriate measure of the support-gradient strength; H3, that large surrogate gain destabilizes the model through coupling to the score-scale degree of freedom; and H4, that \(\neff\to1\) is mainly a late-stage feature of the instability rather than its initial cause.

The main conclusions are:

\begin{itemize}[leftmargin=1.4em]
\item \textbf{H1 is supported.} A score-unit reparameterization leaves the forward computation and training dynamics unchanged while changing \(\chigeo\) by an arbitrary factor. Within a fixed \((K,T)\) row, \(\chigeo\) also has approximately chance-level prospective predictive power (AUC \(\approx0.46\)). A universal numerical threshold in \(\chigeo\) therefore cannot be fundamental.

\item \textbf{H2 is supported as a statement about surrogate gain.} \(\Pisurr=\|L_\kappa D_z\|_F^2\) is dimensionless and reparameterization-invariant. Across the measured states it predicts the realized gradient gain with rank correlation \(0.97\), compared with \(0.52\) for \(\chigeo\). The approximation
\[
\Pisurr\approx\frac{b^2}{4T\delta}
\]
is reasonably accurate in healthy, locally dense boundary regions but can fail badly once the score geometry becomes degenerate.

\item \textbf{H3 is supported, but not in the form of a positive first-order radial force.} At dangerous states, the signed task and surrogate contributions to score scale are generally restoring. The destabilizing contribution is instead a second-order increase in scale caused by large parameter updates. We call this contribution \emph{heating}. It is even under gradient reversal and scales approximately as the square of the learning-rate multiplier. At the most dangerous states it is largely deterministic rather than minibatch-noise driven. Instability begins when heating becomes comparable to or larger than the restoring signed drift.

\item \textbf{H4 is supported with a qualification.} In the strict one-feature limit, \(\neff=1\), the kappa-corrected surrogate vanishes and the measured \(\Pisurr\) is zero. However, neighboring blocks with \(\neff\approx2\) and very large activations can retain enormous surrogate gain and continue to increase the score scale. Post-collapse growth stops when the support gradient is disabled, while broadening the kernel after collapse can make the instability substantially worse.

\item \textbf{The controller should use two variables.} The effective kernel population \(\neff\) is the strongest measured within-row early warning of a burst, while \(\Pisurr\) measures support-surrogate strength. After a block has entered the collapsed regime, the appropriate emergency intervention is to suppress the support gradient rather than increase \(T\).
\end{itemize}

% ============================================================
\section{Definitions}
% ============================================================

All gate quantities are computed per layer from the candidate set consisting of the top \(K+J\) features by score.

The pre-gate feature is \(z_i\), the ranking score is \(s_i=|z_i|\), and the boundary is
\[
b=\max(b_0,s_{(K+1)}).
\]
The kernel is
\[
\kappa_i=\frac{e^{-|s_i-b|/T}}{2T},\qquad Z=\sum_i\kappa_i,
\]
with effective population
\[
\neff=\frac{Z^2}{\sum_i\kappa_i^2}
\]
and local rank-spacing estimate
\[
\delta=\frac{s_{(K-3)}-s_{(K+5)}}{8}.
\]
Small \(\delta\) corresponds to a locally dense score distribution near the support boundary.

To make the backward path explicit, let \(p_i\) denote the surrogate gate variable and let \(y_i=p_i z_i\) be the bottleneck output. The incoming gradient is
\[
u_i=\frac{\partial L}{\partial y_i}.
\]
Since \(\partial y_i/\partial p_i=z_i\),
\[
\frac{\partial L}{\partial p}=D_z u,\qquad D_z=\operatorname{diag}(z).
\]
Define
\[
D_\kappa=\operatorname{diag}(\kappa),\qquad
L_\kappa=D_\kappa-\frac{\kappa\kappa^\top}{Z}.
\]
The kappa-corrected surrogate score gradient is
\[
g_s=L_\kappa D_z u.
\]
Equivalently, with \(w_i=-u_i z_i\), one may write \(g_s=-L_\kappa w\), up to the sign convention used for the reported update direction. Since \(L_\kappa\mathbf 1=0\), the kappa correction removes the common score-translation mode.

The exact surrogate gain is defined as
\[
\boxed{\Pisurr=\|L_\kappa D_z\|_F^2.}
\]
Using the diagonal-minus-rank-one structure,
\[
\Pisurr=
\sum_j z_j^2\kappa_j^2
\left[
\left(1-\frac{\kappa_j}{Z}\right)^2+
\frac{\sum_i\kappa_i^2-\kappa_j^2}{Z^2}
\right].
\]
Because each \(\kappa_j z_j\) is dimensionless, \(\Pisurr\) is dimensionless. In contrast,
\[
\chigeo=\frac{b}{2T\delta}
\]
has units of inverse score. In a locally dense, approximately uniform boundary region,
\[
\Pisurr\approx\frac{b^2}{4T\delta}.
\]
This approximation is tested below rather than assumed generally.

The realized empirical gain is
\[
G_{\mathrm{emp}}=\frac{\|g_s\|^2}{\|u\|^2}.
\]

Three score-scale statistics are monitored on a fixed evaluation batch: \(R_{\mathrm{all}}=\operatorname{RMS}_i(s_i)\), the mean top-\(K\) score \(R_K\), and the boundary \(b\). The logarithmic scale coordinate is \(x=\log R\).

Counterfactual optimizer steps are evaluated from the same saved model and optimizer state using \emph{null} (zero current gradient), \emph{task} (support gradient disabled), \emph{full}, and corresponding antithetic branches with the current gradient reversed. The saved Adam-type state is retained in every branch.

For a gradient \(G\), define the signed, or odd, scale response
\[
\mu(G)=\frac12\mathbb E[\Delta x(+G)-\Delta x(-G)].
\]
The support contribution is \(\mu_{\mathrm{supp}}=\mu_{\mathrm{full}}-\mu_{\mathrm{task}}\). The even component relative to the null branch measures the sign-independent scale increase. At second order, parameter displacements can increase squared activation norms even when their first-order radial contribution changes sign; this contribution is referred to as \emph{heating}.

The raw gradients satisfy
\[
G_{\mathrm{full}}=G_{\mathrm{task}}+G_{\mathrm{supp}},
\]
which was verified with identical random-number states and no stochastic layers. Adam-type optimizer steps are not additive in the same way, so the optimizer branches are measured separately.

% ============================================================
\section{Experiments performed}
% ============================================================

\begin{enumerate}[leftmargin=1.6em]
\item \textbf{Checkpoint library.} Thirteen states were selected: calm \(T=2\) states; stable \(K=32,T=1\) states; a marginal surviving \(K=32,J=224,T=0.5\) run at steps 2000, 4000, and 6000; marginal checkpoints from runs that later collapsed, including \(K=32,J=480,T=0.5\) and \(K=64,J=448,T=1\); and a post-collapse \(K=32,J=480,T=0.1\) state. All checkpoints include full optimizer state. Except for the explicitly post-collapse state, each model is still healthy at the measurement checkpoint. For example, the \(K=64,T=1\) pre-collapse checkpoint is at step 2000, while the original run destabilizes near step 3940.

\item \textbf{One-step drift and heating.} Sixteen batches were evaluated per state. Five counterfactual optimizer branches were measured per batch, with per-layer \(\Delta\log R\) evaluated on a fixed evaluation batch.

\item \textbf{Surrogate gain and radial structure.} For each state and layer the analysis measured
\[
\Pisurr,\quad \frac{b^2}{4T\delta},\quad \chigeo,\quad \neff,\quad G_{\mathrm{emp}},
\]
together with
\[
r_{\mathrm{supp}}=-\frac{s^\top g_s^{(z)}}{s^\top s},\qquad
(s-b\mathbf1)^\top L_\kappa w,\qquad
\operatorname{Cov}_q(s-b,w).
\]
A permutation control used 48 random shuffles of \(w\) among candidates for each token.

\item \textbf{Scale-response curves.} A function-preserving gate-level intervention rescales \(z\to\alpha z\) while dividing the gate output by \(\alpha\). The forward function and ordinary task-gradient path remain unchanged, while the support surrogate sees scores multiplied by \(\alpha\) at fixed \(T\). Measurements were taken at
\[
\alpha\in\{0.5,0.75,1,1.25,1.5,2,3\}
\]
for four reference states.

\item \textbf{Score-unit reparameterization.} The transformation
\[
z\to cz,\qquad s\to cs,\qquad b\to cb,\qquad T\to cT,
\]
with the corresponding output compensation, should leave the gate and parameter gradients invariant while sending \(\chigeo\to\chigeo/c\).

\item \textbf{Learning-rate and batch-size tests.} Learning-rate multipliers \(\{0.5,1,2\}\) and a half-batch condition were tested at three states.

\item \textbf{Counterfactual continuations.} Longer continuations were run from a marginal \(K=96,T=1\) state and from the post-collapse \(T=0.1\) state. Interventions included normal continuation with three seeds, disabling the support gradient immediately or after a boundary-inflation trigger, and increasing \(T\) by a factor of 3 in the marginal state or 30 in the post-collapse state.

\item \textbf{Prospective comparison.} A dataset of 148 checkpoint states from 18 earlier fixed-\(T\) runs was used to test whether a quantity measured at step \(r\) predicts a boundary burst during \((r,r+H]\), both pooled across configurations and within fixed \((K,T)\) rows.
\end{enumerate}

% ============================================================
\section{Results 1: reparameterization test (H1)}
% ============================================================

The transformation
\[
(s,z,b,T)\to(cs,cz,cb,cT)
\]
with the corresponding output compensation leaves the forward computation and parameter gradients invariant.

This invariance is exact in float32 single-gate unit tests and is reproduced at full checkpoints. For \(c=0.5\) and \(c=2\), which are exactly representable binary rescalings, the training loss is unchanged to all reported digits. The largest relative parameter-gradient difference is approximately
\[
2\times10^{-3}\text{--}6\times10^{-3},
\]
consistent with bf16 accumulation differences.

For \(c=3\), some nearly tied scores change ordering in bf16, producing a loss shift of order \(10^{-3}\). This is attributed to finite-precision ranking rather than failure of the underlying invariance.

Under the same transformation,
\[
\chigeo\to\frac{\chigeo}{c}.
\]
For example, in layer 4 the measured values are \(54.1\) at \(c=1\), \(27.0\) at \(c=2\), and \(18.1\) at \(c=3\), while \(\Pisurr\) remains invariant.

Thus functionally equivalent parameterizations can have different values of \(\chigeo\). This rules out a universal numerical instability threshold based directly on \(\chigeo\).

The prospective analysis in Section~\ref{sec:prospective} gives a second limitation: within fixed \((K,T)\) rows, \(\chigeo\) has approximately chance-level predictive power.

\paragraph{H1 conclusion.}
The data support treating \(\chigeo\) as a configuration-level empirical proxy rather than a fundamental order parameter.

% ============================================================
\section{Results 2: surrogate gain \texorpdfstring{$\Pi$}{Pi} (H2)}
% ============================================================

Across 96 layer states (12 checkpoints \(\times\) 8 blocks, excluding the post-collapse state), the exact quantity
\[
\Pisurr=\|L_\kappa D_z\|_F^2
\]
predicts the realized gain \(G_{\mathrm{emp}}=\|g_s\|^2/\|u\|^2\) with rank correlation \(0.97\). The corresponding correlation for \(\chigeo\) is \(0.52\).

The dense-boundary approximation
\[
\Pisurr\approx\frac{b^2}{4T\delta}
\]
is typically accurate within a factor of about two in mid-depth blocks of healthy states. Its accuracy deteriorates when the local score distribution becomes nonuniform or degenerate: discrepancies reach roughly a factor of 10 during late inflation and \(10^4\) after collapse.

The approximation can also overestimate the gain by orders of magnitude when several ranks become nearly tied. In that case \(\delta\to0\) even though the exact kernel geometry and exact \(\Pisurr\) remain finite. This is an important limitation of any quantity based directly on a single local-spacing estimate.

\begin{figure}[h]\centering
\includegraphics[width=.9\textwidth]{figures/sdn\_Pi.png}
\caption{Left: exact surrogate gain \(\Pisurr\) versus the realized gradient gain \(\|g_s\|^2/\|u\|^2\). Right: exact \(\Pisurr\) versus \(b^2/(4T\delta)\). Each point is one checkpoint/block pair. The checkpoint classes describe eventual run fate; only the post-collapse points are measured after collapse. No averaging over blocks is performed.}
\end{figure}

The previous comparison between the \(K\)- and \(T\)-axes is approximately reproduced by the dense-boundary formula. At step 2000,
\[
\frac{\Pi_{\mathrm{approx}}(K=64,T=1)}
{\Pi_{\mathrm{approx}}(K=32,T=1)}
=1.66,
\]
within the predicted range \(1.6\)--\(2.25\). For the exact gain, however, the corresponding ratio is only \(1.09\), because the exact expression uses the full kernel and \(z\)-profile rather than a locally uniform approximation.

\paragraph{H2 conclusion.}
The data support \(\Pisurr\) as the appropriate measure of instantaneous support-surrogate gain. It is dimensionless, reparameterization-invariant, and closely predicts realized gradient amplification. Large \(\Pisurr\), however, is not by itself evidence of instability.

% ============================================================
\section{Results 3: signed drift and second-order heating (H3)}
% ============================================================

\begin{figure}[h]\centering
\includegraphics[width=\textwidth]{figures/sdn\_decomposition.png}
\caption{One-step change in \(\log R_{\mathrm{all}}\), decomposed into signed task drift, signed support-surrogate drift, and even heating, pooled over blocks at thirteen checkpoint states.}
\end{figure}

Three main results emerge.

\begin{enumerate}[leftmargin=1.6em]
\item \textbf{The signed surrogate drift is not the main runaway term.} At the \(K=64,T=1\) checkpoint measured approximately 1940 steps before collapse, the signed surrogate contribution is strongly restoring:
\[
\mu_{\mathrm{supp}}\approx-5.7\times10^{-2}
\]
per step pooled over blocks and approximately \(-0.10\) at the most extreme block. The task contribution is also restoring, \(\mu_{\mathrm{task}}\approx-2.2\times10^{-2}\). A similar pattern appears during late inflation of the unstable \(T=0.5\) run.

\item \textbf{The dominant destabilizing contribution is second-order heating.} At the pre-collapse state, the even contribution reaches \(+9.2\times10^{-2}\) per step pooled and approximately \(+0.18\) at the most affected block, exceeding the restoring signed drift. At calm and stable states it is approximately 30--300 times smaller. The learning-rate sweep shows approximately linear scaling of the odd terms and quadratic scaling of the even term.

\item \textbf{Dangerous-state heating is mostly deterministic.} At the pre-collapse state, halving the batch size changes heating only from \(9.2\times10^{-2}\) to \(9.3\times10^{-2}\), with small across-batch variance. At the stable \(K=32,T=1\) state, the much smaller heating approximately doubles when the batch is halved, indicating a stronger stochastic-gradient contribution there.
\end{enumerate}

\begin{figure}[h]\centering
\includegraphics[width=\textwidth]{figures/sdn\_lr\_scaling.png}
\caption{\label{fig:lr}Signed drift and heating versus learning-rate multiplier at three states, with linear and quadratic reference scalings.}
\end{figure}

The scale-response experiment isolates the signed support contribution. At the calm \(K=64,T=2\) state, the support drift crosses zero near \(\alpha\approx1.4\): it increases scale below that point and becomes restoring above it. At the stable \(K=32,T=1\) state, the signed surrogate contribution remains positive over the measured range but decreases in magnitude with increasing \(\alpha\). The task drift is \(\alpha\)-independent by construction and is measured to be so.

\begin{figure}[h]\centering
\includegraphics[width=\textwidth]{figures/sdn\_Falpha.png}
\caption{\label{fig:falpha}Signed surrogate and task drift under a function-preserving scale intervention \(z\to\alpha z\).}
\end{figure}

The permutation test shows that the relationship between task preference \(w_i\) and rank matters for the signed radial component. The true value differs from the rank-shuffled null by typically \(3\)--\(20\) standard deviations. At the pre-collapse state the alignment is outward relative to the shuffled null (\(+19.8\sigma\)), while late calm states show inward alignment. The absolute signed contribution is nevertheless small compared with heating at dangerous checkpoints.

\paragraph{H3 conclusion.}
Large surrogate gain does not destabilize the model mainly through a positive first-order radial force. The measurements instead support
\[
\text{large support-induced step}
\longrightarrow
\text{second-order scale heating},
\]
opposed by signed task and support restoration. Local instability occurs when heating becomes comparable to or larger than restoration.

% ============================================================
\section{Results 4: counterfactual interventions and collapse dynamics (H4)}
% ============================================================

\begin{figure}[h]\centering
\includegraphics[width=\textwidth]{figures/sdn\_interventions.png}
\caption{Left: continuations from the marginal \(K=96,T=1\) state. Middle: continuations from the collapsed \(T=0.1\) state. Right: one-step drift/heating decomposition along a densely sampled \(K=96,T=0.75\) collapse trajectory.}
\end{figure}

\textbf{Marginal-state intervention.}
Starting from the \(K=96,T=1\) checkpoint, three normal 1500-step continuations show repeated score-scale excursions of one to three orders of magnitude. Most recover, but one reaches terminal divergence.

When the support gradient is disabled, the trajectory becomes calm and score scale decreases. Increasing \(T\) by a factor of three also suppresses the bursts while keeping scale approximately constant, and gives evaluation CE \(1.505\), compared with \(1.552\)--\(1.634\) for the normal branches. Both interventions remain effective when applied only at the first boundary-inflation trigger. Thus, at this checkpoint, the support surrogate is necessary for burst initiation.

\textbf{Post-collapse intervention.}
The collapsed \(T=0.1\) checkpoint begins at scale \(\sim10^6\). Under normal continuation, the scale increases by another factor of approximately 25 over 1200 steps while CE remains near \(5.8\).

Increasing \(T\) by a factor of 30 accelerates the inflation to approximately a factor of 500. Broadening the kernel repopulates candidate regions where some blocks already have
\[
\Pisurr\sim10^{12}.
\]
At the same time, one block is exactly at
\[
\neff=1,\qquad \Pisurr=0,
\]
consistent with \(L_\kappa\to0\) in the strict one-feature limit.

Other near-degenerate blocks still contain more than one significantly weighted feature and have extremely large \(z\), so their surrogate gain remains large.

Disabling the support gradient stops the continued inflation and evaluation CE improves from \(5.93\) to \(4.44\) over 1200 steps. Thus post-collapse inflation is not maintained by the task gradient alone.

\textbf{Dense collapse trajectory.}
In a \(K=96,T=0.75\) run sampled every 100 steps, approximately 200 steps before collapse the measured terms are still moderate:
\[
\text{restoring drift}\approx-1.3\times10^{-2},\qquad
\text{heating}\approx+2.4\times10^{-3}.
\]
Around 60--100 steps before collapse, both grow by roughly two orders of magnitude:
\[
\text{surrogate drift}\approx-0.50,\qquad
\text{heating}\approx+0.37.
\]
The failure therefore involves large competing restoring and destabilizing contributions rather than a simple unopposed outward force. After collapse, both terms become small because much of the kernel is effectively inactive.

\paragraph{H4 conclusion.}
The data support treating \(\neff\to1\) as a late-stage consequence rather than the initial trigger. The exact \(\neff=1\) limit eliminates the corrected support gradient, but nearby blocks with \(\neff\approx2\) and very large activations can retain enormous surrogate gain and continue to drive scale growth.

% ============================================================
\section{Results 5: prospective comparison of candidate indicators}
\label{sec:prospective}
% ============================================================

The candidate quantities were evaluated prospectively on 148 checkpoint states from 18 earlier fixed-\(T\) runs, excluding post-collapse states. For a checkpoint at step \(r\), the target is whether a boundary-inflation burst occurs during \((r,r+H]\).

\begin{center}\small
\begin{tabular}{lcccccc}
\toprule
& \multicolumn{3}{c}{pooled AUC}
& \multicolumn{3}{c}{within-row AUC} \\
\cmidrule(lr){2-4}\cmidrule(lr){5-7}
\(H\)
& \(\chigeo\)
& \(\Pi_{\mathrm{approx}}\)
& \(1/\neff\)
& \(\chigeo\)
& \(\Pi_{\mathrm{approx}}\)
& \(1/\neff\) \\
\midrule
1000 & 0.81 & \textbf{0.88} & 0.77 & 0.46 & 0.63 & \textbf{0.75} \\
2000 & 0.78 & \textbf{0.86} & 0.71 & 0.48 & 0.68 & \textbf{0.75} \\
4000 & 0.80 & \textbf{0.85} & 0.68 & 0.44 & 0.69 & \textbf{0.76} \\
\bottomrule
\end{tabular}
\end{center}

When configurations are pooled, \(\Pi_{\mathrm{approx}}\) gives the best predictive performance. This comparison, however, contains strong between-row differences.

Within fixed \((K,T)\) rows, \(\chigeo\) falls to approximately chance level, with AUC \(0.44\)--\(0.48\). The strongest predictor is low effective population \(1/\neff\), with AUC \(0.75\)--\(0.76\). The approximate surrogate gain remains informative but weaker, with AUC \(0.63\)--\(0.69\).

Only three marginal \((K,T)\) rows contain both positive and negative outcomes, so the within-row values should be regarded as preliminary. The qualitative distinction is nevertheless clear: \(\chigeo\) mostly separates operating regimes, whereas decreasing \(\neff\) contains additional temporal information about an approaching burst.

% ============================================================
\section{Results 5b: surrogate-gain fate map}
% ============================================================

The earlier dose--response analysis was repeated using the surrogate-gain quantities, resolved by block and checkpoint. Only pre-collapse checkpoints are included. The trajectory panel follows the highest-gain block of each run over training, while the dose--response panel compares pre-onset burst rate with the earliest healthy worst-block gain.

The current dataset includes eighteen fixed-\(T\) campaign runs with checkpointed score ladders for \(K=32,64\). A separate sixteen-configuration boundary-bracketing campaign over \(K=32\)--128 was not yet complete.

\begin{figure}[h]\centering
\includegraphics[width=.94\textwidth]{figures/sdn\_pi\_fate\_map.png}
\caption{Exact surrogate gain \(\Pisurr\) versus run fate, block, and training step. The worst-block gain rises toward collapse in failing runs and generally decreases in surviving runs. The surviving rise-and-return trajectory corresponds to the recoverable \(T=0.5\) burst.}
\end{figure}

\begin{figure}[h]\centering
\includegraphics[width=.94\textwidth]{figures/sdn\_piapprox\_fate\_map.png}
\caption{The same analysis using \(b^2/(4T\delta)\). The broad ordering remains, but within-block separation is weaker and near-zero \(\delta\) produces extreme outliers. The exact \(\Pisurr\) is preferable when the kernel tensors are available.}
\end{figure}

% ============================================================
\section{Results 6: drift, heating, and the escape picture}
% ============================================================

The learning-rate and batch-size experiments separate two contributions to the scale dynamics. Signed drift scales approximately linearly with the learning-rate multiplier, while heating scales approximately quadratically.

At the stable \(K=32,T=1\) state, halving the batch size approximately doubles heating, indicating a substantial stochastic-gradient contribution. At the pre-collapse \(K=64,T=1\) state, halving the batch leaves heating essentially unchanged, indicating deterministic second-order overshoot from a large structured update.

The stochastic-looking burst phenomenology is reproduced directly in the marginal continuation experiment: from the same marginal checkpoint, one of three replicas collapses while two undergo bursts and recover.

These observations are consistent with a metastable escape picture, but the current data do not yet support a quantitative reduction to one Kramers-like escape number. At the calm \(K=64,T=2\) state, for example, the restoring scale-response curve crosses zero near \(\alpha\approx1.4\), with local slope
\[
\lambda\approx7\times10^{-3}
\]
per unit \(\log R\) per step. Comparable estimates of restoring slope, excitation strength, and distance to collapse are not available across enough marginal configurations.

The qualitative picture is therefore
\[
\text{restoring scale dynamics}
+\text{second-order excitation}
+\text{kernel-population collapse}.
\]
The task gradient generally supplies restoration, and the support surrogate can also provide restoring signed drift around its scale-dependent operating point. Large support-induced updates simultaneously produce second-order heating. If an excursion becomes sufficiently large, \(\neff\) decreases and the correction enters a near-degenerate regime.

% ============================================================
\section{Proposed causal sequence for a burst}
% ============================================================

Measured facts are marked \textbf{(M)}; remaining interpretation is marked \textbf{(I)}.

\begin{enumerate}[leftmargin=1.6em]
\item In marginal configurations, \(\Pisurr\) is approximately 3--10 times larger than in calm states, and the optimizer step produces heating comparable to or larger than restoring drift in the most affected blocks. \textbf{(M)}

\item This heating produces temporary increases in score scale. The signed task and support contributions generally oppose the inflation at the measured dangerous states, so many excursions recover. \textbf{(M)}

\item During sufficiently large excursions, \(\neff\) decreases. This is observed directly in the \(T=0.1\) collapse trajectory, and low \(\neff\) is the strongest measured within-row predictor of a future burst. \textbf{(M)}

As the population becomes small, the correction is distributed over fewer effective directions, which may increase the gain of the remaining nondegenerate directions. \textbf{(I)}

\item In the exact single-feature limit, \(\neff=1\), the corrected support Jacobian vanishes and \(\Pisurr=0\). At the same time, neighboring blocks with \(\neff\approx2\) and very large \(z\) can have \(\Pisurr\sim10^{12}\). Disabling the support gradient stops post-collapse scale growth, showing that these remaining surrogate contributions are causally important. \textbf{(M)}

\item Once several blocks enter the high-scale regime, scale amplification propagates through depth in the residual-replacement architecture and the trajectory eventually reaches a high-scale state in which most of the kernel is effectively inactive. \textbf{(M)}
\end{enumerate}

% ============================================================
\section{Recommended controller variables}
% ============================================================

The experiments suggest separating two functions that were previously combined in the \(\chigeo\)-based controller.

\textbf{Fast warning: effective kernel population.}
The strongest measured within-row warning signal is \(\neff\). Low \(\neff\) reaches AUC approximately \(0.75\), and population collapse is directly associated with entry into the near-degenerate regime.

The existing controller already contains a related population floor. The present results support keeping such a guard, but the intervention should depend on the stage of the instability. Before collapse, broadening the kernel can restore a healthy distributed correction. After a block has already entered the extreme high-scale regime, increasing \(T\) can repopulate the kernel while \(z\) is enormous and produce very large surrogate gain. In that regime, the measured safe intervention is to suppress the support gradient until the scale recovers.

\textbf{Pressure gauge: exact surrogate gain.}
The natural support-surrogate strength measure is
\[
\Pisurr=\|L_\kappa D_z\|_F^2.
\]
It is dimensionless, reparameterization-invariant, and closely tracks realized support-gradient gain. When exact tensors are unavailable,
\[
\Pi_{\mathrm{approx}}=\frac{b^2}{4T\delta}
\]
is a useful coarse proxy, but it becomes unreliable in degenerate score geometries.

A controller based directly on \(\Pisurr\) is therefore a more principled replacement for the current \(\chigeo\)-based pressure trim. The existing \(\chigeo\) controller works in the tested regime because \(\chigeo\) co-varies with surrogate gain across configurations, but \(\chigeo\) itself has no invariant interpretation.

\textbf{Direct dynamical indicator.}
The quantity most directly related to imminent scale instability is the competition between heating and restoration at the most affected block. Define schematically
\[
\mathcal R_{\mathrm{heat}}
=
\frac{\text{second-order heating}}
{|\text{restoring signed drift}|}.
\]
Measured values are approximately \(0.02\) at calm states, \(0.7\)--\(1.8\) at states whose runs collapse within roughly 2000 steps, and \(\sim2\) about 100 steps before collapse. Thus the direct dynamical criterion is
\[
\boxed{\text{heating becomes comparable to or larger than restoration}.}
\]

The practical forward-only indicators are the pair
\[
(\Pisurr,\neff)
\]
evaluated per block. \(\Pisurr\) measures the squared gain of the surrogate Jacobian and therefore tracks the squared support-gradient amplification relevant to second-order heating. \(\neff\) measures proximity to the near-degenerate regime.

Because \(\Pisurr\) is already a squared gain, one should not generally expect heating to scale as \(\Pisurr^2\). Under otherwise fixed conditions, approximately linear scaling of heating with \(\Pisurr\) is the more natural leading expectation.

% ============================================================
\section{Remaining uncertainties and next experiments}
% ============================================================

\begin{itemize}[leftmargin=1.4em]
\item The destabilizing heating contribution has been identified locally, but transfer of scale between successive blocks has not been measured directly. A block-resolved intervention measuring how scale growth in one block changes the next block within the same optimizer step would test the proposed depth-propagation mechanism.

\item The drift/heating decomposition is based mainly on one-step counterfactuals. Adam's second-moment adaptation changes the response to repeated large gradients, so matched 10--50 step continuations from the same checkpoints are needed to measure the closed-loop attenuation.

\item The within-row prospective comparison contains only three \((K,T)\) rows with both outcomes. The interrupted sixteen-configuration boundary-bracketing campaign would substantially strengthen the comparison between \(\Pisurr\), \(\neff\), and other candidate warning variables.

\item A controller targeting exact \(\Pisurr\) has not yet been tested. The present analysis predicts that it should be more portable than the \(\chigeo\)-based controller under changes of score units and parameterization.

\item The metastable-escape interpretation remains qualitative. A quantitative escape parameter would require matched estimates of quantities analogous to the restoring slope, excitation strength, and effective collapse boundary, e.g.\ \((\lambda,D,x_c)\), across a larger number of marginal configurations.
\end{itemize}

\end{document}